%! AP Statistics — Unit 8, Lesson 8.5 — Homework
\documentclass[10pt]{article}
\usepackage{apstats-article}
\usepackage{apstats-boxes}
\usepackage{grading}
\gradingassignment{apstats-u08-l05-hw}
\gradingcourse{advanced_statistics}
\newcommand{\TallMath}[1]{$\displaystyle #1\rule[-1.4em]{0pt}{3.2em}$}

\begin{document}

\pageheader{Unit 8, Lesson 8.5}{Homework}
{\small\itshape Write each final answer in its box. Only what is inside a box is graded.}\par

\vspace{0.1in}

\begin{enumerate}

  \item For each scenario below, \textbf{(i)} identify whether the appropriate test is a
        chi-square test for homogeneity or independence, \textbf{(ii)} explain your reasoning
        based on the study design, and \textbf{(iii)} write $H_0$ and $H_a$ in context.

        \begin{enumerate}[label=(\alph*)]

          \item Three brands of coffee (Brand X, Y, Z) each recruited a random sample of
                50 coffee tasters who rated their experience as Prefer, Neutral, or Dislike.

                Test type: \makebox[4cm][r]{\smash{\gradedblank[id=hw1a1, type=short, points=1, answer=homogeneity, check=text]{2.2cm}{0.2in}}}

                Reason: \gradedblank[id=hw1a2, type=short, points=1, answer={separate samples, one variable}, check=text, display={\bfseries Separate samples, one variable.}]{4.5in}{0.3in}

                $H_0$: \gradedblank[id=hw1a3, type=short, points=1, answer={same distribution of ratings for X, Y, Z}, check=text, display={\bfseries Same for X, Y, Z.}]{4.5in}{0.3in}

                \vspace{3pt}
                $H_a$: \gradedblank[id=hw1a4, type=short, points=1, answer={not the same for all three brands}, check=text, display={\bfseries Not all the same.}]{4.5in}{0.3in}
                \vspace{3pt}

          \item A simple random sample of 400 shoppers is asked their age group (18--30,
                31--50, or 51+) and their preferred shopping method (In-store, Online, or Both).

                Test type: \makebox[4cm][r]{\smash{\gradedblank[id=hw1b1, type=short, points=1, answer=independence, check=text]{2.2cm}{0.2in}}}

                Reason: \gradedblank[id=hw1b2, type=short, points=1, answer={one sample, two variables}, check=text, display={\bfseries One sample, two variables.}]{4.5in}{0.3in}

                $H_0$: \gradedblank[id=hw1b3, type=short, points=1, answer={no association}, check=text, display={\bfseries No association.}]{4.5in}{0.3in}

                \vspace{3pt}
                $H_a$: \gradedblank[id=hw1b4, type=short, points=1, answer={association}, check=text, display={\bfseries Association.}]{4.5in}{0.3in}
                \vspace{3pt}

          \item A simple random sample of 200 high school athletes is classified by sport
                (Football, Basketball, or Soccer) and GPA level (3.0 and above, or Below 3.0).

                Test type: \makebox[4cm][r]{\smash{\gradedblank[id=hw1c1, type=short, points=1, answer=independence, check=text]{2.2cm}{0.2in}}}

                Reason: \gradedblank[id=hw1c2, type=short, points=1, answer={one sample, two variables}, check=text, display={\bfseries One sample, two variables.}]{4.5in}{0.3in}

                $H_0$: \gradedblank[id=hw1c3, type=short, points=1, answer={no association}, check=text, display={\bfseries No association.}]{4.5in}{0.3in}

                \vspace{3pt}
                $H_a$: \gradedblank[id=hw1c4, type=short, points=1, answer={association}, check=text, display={\bfseries Association.}]{4.5in}{0.3in}
                \vspace{3pt}

        \end{enumerate}

  \newpage
  \item A student says scenarios (b) and (c) in Problem 1 both involve two categorical
        variables, so both must be tests for independence. Is the student correct?
        Explain your reasoning.

        \gradedhide{\gradedhere{\gradedblank[id=hw2, type=rubric, points=2, display={Conclusion right (both are independence), reasoning wrong: two variables do not decide it --- the design does (one sample, two variables = independence; separate samples, one variable = homogeneity).}]{\linewidth+6pt}{\gradedlinesheight{4}}}\writelines{9}}

  \item A researcher randomly selects 50 students from each of 3 schools
        (School P, Q, R) and asks each student which after-school program they prefer
        (Academic Support, Sports, or Arts). Total $n = 150$.

        \begin{enumerate}[label=(\alph*)]
          \item Identify the test type and write the hypotheses.

                \gradedhide{\gradedhere{\gradedblank[id=hw3a, type=rubric, points=3, display={Homogeneity. $H_0$: program-preference distribution is the same for Schools P, Q, R. $H_a$: not the same for all three.}]{\linewidth+6pt}{\gradedlinesheight{3}}}\writelines{9}}

          \newpage
          \item Write out the two conditions for this chi-square test.
                Can either condition be verified without seeing the actual data table?
                Explain.

                \gradedhide{\gradedhere{\gradedblank[id=hw3b, type=rubric, points=3, display={Random/10\%: 50 random students per school; need each school $\ge 500$ students (needs school sizes). Large counts: all expected counts $\ge 5$ --- needs the table. Neither can be verified yet.}]{\linewidth+6pt}{\gradedlinesheight{4}}}\writelines{9}}
        \end{enumerate}

  \item \textbf{AP-Style Multiple Choice.} Which of the following scenarios calls for
        a chi-square test for \textbf{homogeneity}? Write the letter of the single best answer in its Letter box.

        \begin{enumerate}[label=(\Alph*)]
          \item A random sample of 500 adults is asked their employment status
                (Employed/Unemployed) and their level of financial stress (Low/Medium/High).
          \item Researchers randomly select 80 men and 80 women and record each person's
                preferred exercise type (Cardio/Strength/Flexibility).
          \item A random sample of 200 teenagers is asked their screen-time category
                ($<$2 hrs/$2$--$4$ hrs/$>$4 hrs/day) and their sleep quality (Good/Poor).
          \item A SRS of 300 patients from one hospital is classified by diagnosis
                (Diabetes/Hypertension/Neither) and insurance type (Private/Public/None).
                      \par\nobreak\rule{0pt}{0.45in}\hfill\gradedright{\textit{Letter:}~\gradedblank[id=hw4, type=choice, points=1, answer=B]{0.4in}{0.2in}}
        \end{enumerate}

\end{enumerate}

\begin{reflectionbox}
In your own words, what is the single most important question to ask when deciding whether
to use a chi-square test for homogeneity or independence? Write 1--2 sentences.
\writelines{3}
\end{reflectionbox}

\end{document}
